\clearpage
\section{The polynomial theorem and the missing integer analogue}\label{sec:mason}
\subsection{Roots play the role of distinct primes}
For a nonzero polynomial $F$ over an algebraically closed field, let $\operatorname{rad}(F)$ be the product of its distinct monic linear factors. Its degree counts distinct roots. Write $n_0(F)=\deg\operatorname{rad}(F)$.

The analogy is direct. Prime powers $p^e$ become repeated roots $(X-\alpha)^e$. A derivative lowers the multiplicity by one. The divisor left over after that operation is $F/\operatorname{rad}(F)$, just as $n/\rad(n)$ is the repeated-power part of an integer.

\begin{theorem}[Mason--Stothers in characteristic zero]\label{thm:mason}
Let $k$ be an algebraically closed field of characteristic zero. Suppose $A,B,C\in k[X]$ are nonzero, pairwise coprime, satisfy $A+B=C$, and are not all constant. Then
\[
 \max\{\deg A,\deg B,\deg C\}\le n_0(ABC)-1.
\]
\end{theorem}
This is the polynomial $abc$ theorem. Snyder gives a short primary proof in a form that also makes the derivative hypothesis explicit~\cite{Snyder2000}. The proof below spells out the cancellation that will motivate the arithmetic Wronskian.

\begin{proof}
Set $W=A'B-AB'$. It is nonzero. Indeed, if it vanished, then $(A/B)'=0$, and characteristic zero would force the rational function $A/B$ to be constant. Coprimality would then force both $A$ and $B$ to be constant, and so also $C$.

If $\alpha$ is a root of $A$ of multiplicity $e$, then $A'$ is divisible by $(X-\alpha)^{e-1}$ and $AB'$ by $(X-\alpha)^e$. Thus $(X-\alpha)^{e-1}$ divides $W$. The same argument applies to roots of $B$.

Since $C=A+B$,
\[
 W=A'C-AC'.
\]
At a root of $C$ of multiplicity $e$, this expression is divisible by $(X-\alpha)^{e-1}$. The three root sets are disjoint. Therefore
\[
 \frac{ABC}{\operatorname{rad}(ABC)}\mid W.
\]
Taking degrees gives
\[
 \deg A+\deg B+\deg C-n_0(ABC)\le\deg W\le\deg A+\deg B-1.
\]
Hence $\deg C\le n_0(ABC)-1$. To bound $A$, apply the same argument to the relation $C+(-B)=A$. To bound $B$, use $C+(-A)=B$. Coprimality and the radical root count are unchanged by these permutations and signs.
\end{proof}

\subsection{A sharp example}
Take $A=X^n$, $B=1-X^n$, $C=1$ for $n\ge1$. The root sets are disjoint. There is one distinct root of $A$ and $n$ distinct roots of $B$, so $n_0(ABC)=n+1$. The maximal degree is $n$, attaining the theorem's bound. In this example $W=nX^{n-1}$, and the repeated-root divisor is exactly $X^{n-1}$ up to a nonzero scalar.

The proof needs an actual nonzero derivative. In positive characteristic, $X^p+(1-X^p)=1$ has every derivative zero, so the characteristic-zero statement cannot simply be transferred. The correct positive-characteristic forms include a condition excluding common Frobenius powers.

\subsection{Where the integer argument stops}
An ordinary integer has no distinguished derivative with both an exact Leibniz rule and an automatic relation $D(a)+D(b)=D(c)$. Arithmetic derivatives provide a family of Leibniz maps. Additivity at one triple then imposes a lattice constraint on the prime-coordinate values.

The polynomial theorem gets two benefits simultaneously. Repeated roots force a large divisor of $W$, while degree controls $W$ by a sharp expression one less than the input degree sum. For integers, a prime-coordinate Wronskian also has a large divisor. Its absolute size depends on the selected derivative packet. Choosing an additive transverse packet of uniformly small height is a separate arithmetic problem.

The subsequent prime-packet construction makes this distinction exact. Its local divisibility is the analogue of the displayed polynomial divisor. Its packet norm, integral synthesis and conductor comparisons are the quantities for which the polynomial proof has no automatic integer replacement.

\subsection{A lesson for higher-order methods}
High contact has value only when the auxiliary object remains nonzero and its global size is controlled. A determinant can raise contact through actual Taylor cancellation. Merely renaming a coordinate with a larger weight does not change its first nonzero Taylor coefficient.

The Hankel determinant later in this piece addresses the nonzero-contact side with an exact coefficient formula and real positivity. Its rational truncation provides a rigorous test of the other side: after projective denominator clearing, the desired output-prime contact disappears. The two obligations must therefore be proved together for a global witness.
